% Ngân hàng tự sinh (AI)
\begin{bt}%[0H9N1-1]
% nguon: AI soạn tương tự câu mẫu 959cabbe
% mau: 959cabbe
Trên trục tọa độ $\left(O;\overrightarrow{i}\right)$ cho 2 điểm $A$, $B$ có tọa độ lần lượt là $-5; 3$. Tìm tọa độ của véc-tơ $\overrightarrow{AB}$.
	\loigiai{
		Ta có  $\overrightarrow{AB}=3+5=8$.\\	Vậy toạ độ $\overrightarrow{AB} $ là $8$.
	}
\end{bt}

\begin{bt}%[0H9N1-2]
% nguon: AI soạn tương tự câu mẫu ae5e3c24
% mau: ae5e3c24
	Trong mặt phẳng tọa độ $Oxy$, cho $2$ véc-tơ $\overrightarrow{a}=(5;3)$, $\overrightarrow{b}=(-6;10)$. Tính tích vô hướng $\overrightarrow{a}\cdot \overrightarrow{b}$ và chứng tỏ $\overrightarrow{a}\perp\overrightarrow{b}$.
	\loigiai{
		Ta có $\overrightarrow{a}\cdot \overrightarrow{b} = 5\cdot (-6) + 3\cdot 10 =0$.\\
		Suy ra $\overrightarrow{a}\perp\overrightarrow{b}$.
	}
\end{bt}

\begin{bt}%[0H9N1-3]
% nguon: AI soạn tương tự câu mẫu 04017960
% mau: 04017960
	Trong mặt phẳng tọa độ $Oxy$, cho hai điểm $P(2; -5)$ và $Q(6; 3)$. Tìm tọa độ $\overrightarrow{PQ}$.
	\loigiai{
		Ta có $\overrightarrow{PQ}=(6-2; 3-(-5))=(4; 8)$.
	}
\end{bt}

\begin{bt}%[0H9H1-4]
% nguon: AI soạn tương tự câu mẫu 2f967af4
% mau: 2f967af4
Trong mặt phẳng $Oxy$, cho $A\left(-2; 1\right), B\left(1; 4\right), C\left(3; -2\right)$. Gọi $M$ là điểm thuộc $BC$ sao cho $\overrightarrow{MB}=-2 \overrightarrow{BC}$ và $G$ là trọng tâm của tam giác $ABC$.
\begin{enumerate}
    \item Phân tích $\overrightarrow{AM}$ theo hai véc-tơ $\overrightarrow{AB}$ và $\overrightarrow{AC}$.
    \item Phân tích $\overrightarrow{GM}$ theo hai véc-tơ $\overrightarrow{AB}$ và $\overrightarrow{AC}$.
\end{enumerate}
\loigiai{
\begin{enumerate}
    \item Gọi $M(x;y) \Rightarrow \overrightarrow{MB}=(1-x;4-y); \overrightarrow{BC}=(2;-6); \overrightarrow{AB}=(3;3); \overrightarrow{AC}=(5;-3)$.\\
    Mà $\overrightarrow{MB}=-2 \overrightarrow{BC} \Leftrightarrow \heva{&1-x=-4\\ &4-y=12} \Leftrightarrow \heva{&x=5\\ &y=-8} \Rightarrow M(5;-8) \Rightarrow \overrightarrow{AM}=(7;-9)$.\\
    Giả sử $\overrightarrow{AM} = m\overrightarrow{AB}+n\overrightarrow{AC} \Leftrightarrow \heva{&7=3m+5n\\ &-9=3m-3n} \Leftrightarrow \heva{&m= -1\\&n=2.}$\\
    Vậy $\overrightarrow{AM} = -\overrightarrow{AB}+2\overrightarrow{AC}$.
    \item Ta có tọa độ trọng tâm $G\left(\dfrac{2}{3};1\right) \Rightarrow \overrightarrow{GM}=\left(\dfrac{13}{3};-9\right)$.\\
    Giả sử $\overrightarrow{GM} = h\overrightarrow{AB}+k\overrightarrow{AC} \Leftrightarrow \heva{&\dfrac{13}{3}=3h+5k\\ &-9=3h-3k} \Leftrightarrow \heva{&h= -\dfrac{4}{3}\\ &k= \dfrac{5}{3}.}$\\
    Vậy $\overrightarrow{GM} = -\dfrac{4}{3}\overrightarrow{AB}+\dfrac{5}{3}\overrightarrow{AC}$.
\end{enumerate}
\textit{\textbf{Chú ý:} Chúng ta hoàn toàn có thể sử dụng cách biến đổi véc-tơ để thực hiện việc biểu diễn $\overrightarrow{AM},\overrightarrow{GM}$ theo hai véc-tơ $\overrightarrow{AB}$ và $\overrightarrow{AC}$.}
}
\end{bt}

\begin{bt}%[0H9N1-5]
% nguon: AI soạn tương tự câu mẫu 42212c6c
% mau: 42212c6c
	Véc-tơ $\overrightarrow{a}=(3;0)$ được phân tích theo hai véc-tơ đơn vị $\left(\overrightarrow{i};\overrightarrow{j}\right)$ như thế nào?
	\loigiai{
		Ta có $\overrightarrow{a}=(3;0)\Rightarrow\overrightarrow{a}=3\overrightarrow{i}+0\overrightarrow{j}=3\overrightarrow{i}$.\\
	}
\end{bt}

\begin{bt}%[0H9H1-6]
% nguon: AI soạn tương tự câu mẫu b279dbb9
% mau: b279dbb9
	Cho $\Delta ABC$ có $A(2;7)$, $B(6;-3)$, $C(-2;1)$. Tọa độ điểm $M(a;b)$ thuộc đoạn $BC$ sao cho $S_{\Delta MAB}=3S_{\Delta MAC}$. Tính giá trị của biểu thức $a+b$.
	\loigiai{
		Kẻ $AH \perp BC$. $S_{\Delta MAB}=3S_{\Delta MAC}\Leftrightarrow \dfrac{1}{2}AH\cdot MB=3\cdot \dfrac{1}{2}AH\cdot MC\Leftrightarrow MB=3MC$.\\
		Mà $\overrightarrow{MB}$ và $\overrightarrow{MC}$ ngược hướng $\Rightarrow \overrightarrow{MB}=-3\overrightarrow{MC}$\\
		$\Leftrightarrow \heva{&x_B-x_M=-3(x_C-x_M)\\&y_B-y_M=-3(y_C-y_M)}\Leftrightarrow \heva{&6-x_M=-3(-2-x_M)\\&-3-y_M=-3(1-y_M)}$\\
		$\Leftrightarrow \heva{&6-x_M=6+3x_M\\&-3-y_M=-3+3y_M}\Leftrightarrow \heva{&4x_M=0\\&4y_M=0}\Rightarrow M(0;0)$.\\
		Suy ra $a+b=0$.
	}
\end{bt}

\begin{bt}%[0H9N2-1]
% nguon: AI soạn tương tự câu mẫu 2dc15a56
% mau: 2dc15a56
	Cho hai véc-tơ $\overrightarrow{a}=(2;1),\overrightarrow{b}=(0;4)$.
	\begin{enumerate}
		\item Tìm toạ độ của véc-tơ $3\overrightarrow{a}+2\overrightarrow{b}$.
		\item Tính các tích vô hướng: $\overrightarrow{a} \cdot \overrightarrow{b},(2\overrightarrow{a}) \cdot (3\overrightarrow{b})$.
	\end{enumerate}
	\loigiai{
		\begin{enumerate}
			\item Ta có: $3\overrightarrow{a}+2\overrightarrow{b}=(3 \cdot 2+2 \cdot 0;3 \cdot 1+2 \cdot 4)=(6;11)$.\\
			\item Ta có: $\overrightarrow{a} \cdot \overrightarrow{b}=2 \cdot 0+1 \cdot 4=4$\\
			$2\overrightarrow{a}=(4;2)$ và $3\overrightarrow{b}=(0;12)$ nên $(2\overrightarrow{a}) \cdot (3\overrightarrow{b})=4 \cdot 0+2 \cdot 12=24$.
		\end{enumerate}
	}
\end{bt}

\begin{bt}%[0H9N2-2]
% nguon: AI soạn tương tự câu mẫu 34f6d564
% mau: 34f6d564
	Cho $\overrightarrow{a}=(3; 2)$ và $\overrightarrow{b}=(-1;5)$. Tính $\cos \alpha$, với $\alpha =\left(\overrightarrow{a}, \overrightarrow{b}\right)$.
	\loigiai{
		Ta có $\cos \alpha =\dfrac{\overrightarrow{a}\cdot \overrightarrow{b}}{\left|\overrightarrow{a}\right|\cdot \left|\overrightarrow{b}\right|} =\dfrac{3\cdot (-1)+2\cdot 5}{\sqrt{3^2+2^2}\cdot \sqrt{(-1)^2+5^2}}=\dfrac{7}{\sqrt{338}}$.\\
	}
\end{bt}

\begin{bt}%[0H9H2-3]
% nguon: AI soạn tương tự câu mẫu a3548811
% mau: a3548811
	Trong mặt phẳng tọa độ $Oxy$, cho tam giác $ABC$ có $A(2;-1)$, $B(3;3)$ và $C(9;1)$. Tính chu vi tam giác $ABC$.
	\loigiai{
		Ta có $\heva{&\overrightarrow{AB}=(1;4)\\&\overrightarrow{AC}=(7;2)\\&\overrightarrow{BC}=(6;-2)} \Rightarrow \heva{&AB=\sqrt{1^2+4^2}=\sqrt{17}\\&AC=\sqrt{7^2+2^2}=\sqrt{53}\\&BC=\sqrt{6^2+(-2)^2}=2\sqrt{10}.}$\\
		Chu vi tam giác $ABC$ là $AB+AC+BC=\sqrt{17}+\sqrt{53}+2\sqrt{10}$.
	}
\end{bt}

\begin{bt}%[0H9N2-4]
% nguon: AI soạn tương tự câu mẫu 0e8eb5bf
% mau: 0e8eb5bf
Trong mặt phẳng tọa độ $Oxy$, cho điểm $A(5;-2)$, đường thẳng $\Delta_1\colon x-2y+7=0$ và đường thẳng $\Delta_2\colon 3x-y-50=0$.
	\begin{enumerate}
		\item Tính góc giữa đường thẳng $\Delta_1$ và $\Delta_2$.
		\item Tính khoảng cách từ điểm $A$ đến đường thẳng $\Delta_2$.
	\end{enumerate}
	\loigiai{
		\begin{enumerate}
			\item $\Delta_1$ có véc-tơ pháp tuyến $\overrightarrow{n}_1=(1;-2)$ và $\Delta_2$ có véc-tơ pháp tuyến $\overrightarrow{n}_2=(3;-1)$.\\
			$\cos(\Delta_1;\Delta_2)=\left|\cos(\overrightarrow{n}_1;\overrightarrow{n}_2)\right|=\dfrac{|1\cdot3+(-2)\cdot(-1)|}{\sqrt{1^2+(-2)^2}\cdot\sqrt{3^2+(-1)^2}}=\dfrac{5}{\sqrt{5}\cdot\sqrt{10}}=\dfrac{\sqrt{2}}{2}$.\\
			Suy ra $(\Delta_1;\Delta_2)=45^\circ$.
			\item $\mathrm{d}(A,\Delta_2)=\dfrac{|3\cdot5-(-2)-50|}{\sqrt{3^2+(-1)^2}}=\dfrac{|15+2-50|}{\sqrt{10}}=\dfrac{33}{\sqrt{10}}=\dfrac{33\sqrt{10}}{10}$.
		\end{enumerate}
	}
\end{bt}
